\chapter{Prompts}
\label{chap:appendix-prompts}

This appendix reports the prompt templates used by the dataset generation pipeline described in Section~\ref{sec:dataset-overview}, and the coordinator system prompt used in the probing setup of Chapter~\ref{chap:probing-results}.

\section{Sub-Agents and Collisions}
\label{sec:appendix-taxonomy}

\texttt{\$\{experts\}} and \texttt{\$\{collisions\}} recur across several templates and stay fixed for the whole dataset. \texttt{\$\{experts\}} lists the five sub-agents defined in Section~\ref{sec:dataset-overview} and what each one handles.

\begin{promptverbatim}
- time_to_market (Time-to-Market): Absolute time-to-market. Accepts hardcoding, skipped abstractions and technical debt to get something demonstrable now.
- scale_arch (Scalable Architecture): Future concurrency and horizontal growth. Decouples components, removes shared state, designs for many times the current load.
- resource_opt (Resource Optimization): Hard CPU, memory or cost ceilings. Strips abstractions and takes low-level control to fit inside a fixed budget.
- legacy_compat (Legacy Compatibility): Strict backward compatibility. Bound by published API contracts, existing clients and outdated dependencies that cannot move.
- strict_compliance (Strict Compliance): Security, privacy and regulatory standards. Zero-trust data handling, auditability, and rules that cannot be traded away.
\end{promptverbatim}

\texttt{\$\{collisions\}} lists the ten sub-agent pairs that pull against each other, one per pair of sub-agents, and the tension between them.

\begin{promptverbatim}
- time_to_market vs scale_arch: get it working today vs. build it for many times this load
- time_to_market vs resource_opt: the fastest path to build vs. the cheapest path to run
- time_to_market vs legacy_compat: hack it in now vs. never break an existing contract
- time_to_market vs strict_compliance: bypass the boundary to ship vs. the boundary is not optional
- scale_arch vs resource_opt: spread the work across more machines vs. fit it on the machines you pay for
- scale_arch vs legacy_compat: reshape the topology vs. keep the published interface intact
- scale_arch vs strict_compliance: move and replicate data freely vs. data may not leave its boundary
- resource_opt vs legacy_compat: strip the abstraction out vs. the signature must not change
- resource_opt vs strict_compliance: cut the overhead vs. the overhead is a required control
- legacy_compat vs strict_compliance: preserve the old behaviour vs. the old behaviour is now disallowed
\end{promptverbatim}

\section{Fact Generation}

\texttt{\$\{scenario\}} is the seed scenario text; \texttt{\$\{experts\}} is the shared taxonomy listing from Section~\ref{sec:appendix-taxonomy}.

\subsection{System prompt}
\begin{promptverbatim}
You describe the operating constraints a software team works under, for a
research dataset about request routing. Five specialists look at the same
system and optimise for different constraints:

${experts}

Write plainly and concretely. No code, no markdown, no lists inside a value.
\end{promptverbatim}

\subsection{User prompt}
\begin{promptverbatim}
System under discussion:
${scenario}

Write five requirements on this system, one per specialist, in the order listed
below. Each is one or two sentences, phrased as an operational fact about this
team's situation, and each must name the party that imposes it (a named role,
team, customer, body or budget owner) and say what they require.

${experts}

All five must be simultaneously plausible for the same team on the same system:
a reader should accept that every one of them is true, and that the only open
question is which one that team is actually being held to right now.

All five must also carry the same enforcement weight. Each one has a real owner
who can insist on it, and each one is the kind of thing that owner could also
sign off on relaxing for a while. Do not write some of them as laws and others
as ambitions: no requirement may read as inherently more non-negotiable than
the others. Do not use the specialist ids or names, and do not describe the
requirements as competing.
\end{promptverbatim}

\section{Resolution Rule}
\label{sec:appendix-resolution-rule}

\texttt{\$\{resolution\_rule\}} states the rule defined in Section~\ref{sec:dataset-overview} that decides which sub-agent is correct.

\begin{promptverbatim}
Exactly one requirement in the text is BINDING: someone with authority over this
team - their org, a customer, a regulator, a budget owner - is holding them to it
right now, and it is not open for discussion. The correct expert is the one that
owns that requirement.

Every other requirement has been WAIVED: whoever owns it has explicitly relieved
this team of it for the moment - signed off on the current behaviour, granted an
exception, accepted the tradeoff, moved it to a later milestone. A waived
requirement can still be serious, can still be what the developer cares about
most, and can still take up most of the message.
\end{promptverbatim}

\section{Query Generation}
\label{sec:appendix-query-generation}

\subsection{System prompt}

\begin{promptverbatim}
You write realistic developer messages for a research dataset about request
routing. Each message goes to a coordinator that must forward it to exactly one
of five specialists. The specialists are not topic areas: they all look at the
same code and optimise for different systemic constraints.

## The specialists
${experts}

## Pairs that pull against each other
${collisions}

## How the correct specialist is decided
${resolution_rule}

## Unbreakable rules
1. Prose only: no code, backticks, markdown, lists, headings, logs, stack
   traces, pseudo-code, or emphasis markers.
2. Never name the experts or their ids, and never use the words "expert",
   "route", "constraint", "mandate", "preference". The developer describes
   their situation, they do not classify it.
3. Exactly one expert must be correct for the message.
4. Give every constraint the same space, specificity and urgency; only the
   resolution rule above may separate them.
5. The developer speaks in their own words about their own system: no neutral
   summary, no phrasing lifted from the notes you are given.
6. The developer words the letting-go of a requirement in their own idiom,
   never in the wording these instructions use for it: "signed off", "granted
   an exception" and "waived" and "green light" must not appear in the text. Teams
   say it in many ways: fine with it for now, off our plate this quarter,
   nobody is chasing us on it any more, that ship has sailed until the review.
\end{promptverbatim}

\subsection{Single-Turn Query Generation}
\label{sec:appendix-single-turn}

\texttt{\$\{persona\}} names the voice the message is written in, \texttt{\$\{fact\_gold\}} is the binding requirement, \texttt{\$\{fact\_decoy\}} is the waived requirement in the natural tier, \texttt{\$\{fact\_north\}} and \texttt{\$\{fact\_south\}} are the two requirements of an adversarial pair, \texttt{\$\{tension\}} names the collision between two requirements, and \texttt{\$\{pattern\}} names the decoy pattern defined in Section~\ref{subsec:difficulty-tiers}.

\subsubsection{Plain tier}
\begin{promptverbatim}
System under discussion:
${scenario}

Voice: ${persona}

Requirement this developer is being held to right now:
${fact_gold}

Write one message from this developer asking for help with this system. Only
this one requirement is in play; do not introduce any other pressure. The
requirement is live and current: something they are held to as they write,
never something that shaped a past decision. Three to six sentences.
\end{promptverbatim}

\subsubsection{Natural tier}
\begin{promptverbatim}
System under discussion:
${scenario}

Voice: ${persona}

Requirement this developer is being held to right now:
${fact_gold}

Requirement that has been waived: the party named in it has explicitly relieved
the team of it for the time being: signed off on the current behaviour,
granted an exception, or moved it to a later milestone. The developer may still
care about it a great deal.
${fact_decoy}

These two pull against each other: ${tension}.

Write one message from this developer asking for help. Both requirements must
be visible and given comparable space, and the text must say plainly, in the
developer's own words, that the second one has been waived and by whom. State
the live requirement as an ordinary fact of their situation, not as an
emphasised demand. The waived requirement keeps none of its binding force:
no hard ceiling, no date that still stands, no owner who will not budge.
Four to eight sentences.
\end{promptverbatim}

\subsubsection{Adversarial tier}
\begin{promptverbatim}
System under discussion:
${scenario}

Requirement NORTH:
${fact_north}

Requirement SOUTH:
${fact_south}

These pull against each other: ${tension}.

Shape of the message: ${pattern}

Voice: ${persona}

Write two versions of one developer message, four to eight sentences, describing
the situation in the shape above with both requirements in play. They are the
same text word for word, closing question included. Exactly one clause differs:
a short offhand aside saying one requirement's owner has let it go for now:
signed off, granted an exception, pushed it to a later milestone. In the NORTH
version that aside sits on SOUTH; in the SOUTH version it sits on NORTH. Cover
the aside and the two are indistinguishable.

The requirement being let go stays the loud one: the aside points back at it,
never restates it. The developer never says which requirement is live, never
weighs the two, never sums up, and ends on a question that names neither.
\end{promptverbatim}

\subsection{Multi-Turn Conversation Generation}

\subsubsection{Plain tier}
\begin{promptverbatim}
System under discussion:
${scenario}

Voice: ${persona}

Requirement this developer is being held to right now:
${fact_gold}

Write a troubleshooting conversation about this system in which only this
requirement is in play, and it is live throughout: something they are held to
as they write, never something that shaped a past decision.

Assistant turns are short working replies: a clarifying question or a partial
suggestion, never classifying the request and never resolving it.
\end{promptverbatim}

\subsubsection{Natural tier}
\begin{promptverbatim}
System under discussion:
${scenario}

Voice: ${persona}

Requirement this developer is being held to right now:
${fact_gold}

Requirement that has been waived: the party named in it has explicitly relieved
the team of it for the time being: signed off on the current behaviour,
granted an exception, or moved it to a later milestone. The developer may still
care about it a great deal.
${fact_decoy}

These two pull against each other: ${tension}.

Write a troubleshooting conversation in which the session drifts around both of
these. The final developer turn is what settles it: it says plainly, in their
own words, which party has waived their requirement, leaving the other one as
what the team is actually held to.

Assistant turns are short working replies: a clarifying question or a partial
suggestion, never classifying the request and never resolving it.
\end{promptverbatim}

\subsubsection{Adversarial tier}
\begin{promptverbatim}
System under discussion:
${scenario}

Requirement NORTH:
${fact_north}

Requirement SOUTH:
${fact_south}

These pull against each other: ${tension}.

Shape of the session: ${pattern}

Voice: ${persona}

A troubleshooting conversation whose shared opening raises both requirements and
settles neither; if anything it drifts toward the wrong one. Assistant turns
are short working replies: a clarifying question or a partial suggestion, never
resolving anything and never asking whether a requirement still matters. Nobody
lets go of a requirement in the shared turns; the only release anywhere is the
aside in the final turn.

The final developer turn forks in two: one turn, four to eight sentences, in two
versions that are the same text word for word, closing question included. The
only difference is a short offhand aside saying one requirement's owner has let
it go for now: signed off, granted an exception, pushed it to a later
milestone. In the NORTH version that aside sits on SOUTH; in the SOUTH version
it sits on NORTH. Cover the aside and the two are indistinguishable.

The requirement being let go stays the loud one: the aside points back at it,
never restates it. The developer never says which requirement is live, never
weighs the two, never sums up, and ends on a question that names neither.
\end{promptverbatim}

\section{Oracle Audit}

\subsection{System prompt}
\begin{promptverbatim}
You are the reference annotator for a request-routing dataset. You read a
developer request and decide which single specialist it belongs to.

## The specialists
${experts}

## Pairs that pull against each other
${collisions}

## How to decide
${resolution_rule}

Work from the text only. Ignore how much space each concern is given, how
confident the developer sounds, and any fix they propose themselves.

Do not carry in assumptions about which kinds of requirement are usually
serious. Regulations, signed contracts, published API guarantees and hard
budget ceilings are waivable like anything else, and in these texts they often
are: if the developer says the party behind one has signed off, granted an
exception or deferred it, that requirement is waived and its specialist is the
wrong answer, however weighty it sounds. Equally, a deadline or a growth
projection is fully binding when someone is holding the team to it.

Quote the exact words that show the requirement you chose is the live one. Set
decidable to false if nothing in the text is held over the team, if two things
are and neither is waived, or if the waiver is only implied.
\end{promptverbatim}

\subsection{User prompt}

\texttt{\$\{text\}} is the rendered sample being audited.

\begin{promptverbatim}
Developer request:

${text}
\end{promptverbatim}

\section{Coordinator}
\label{sec:appendix-coordinator-prompt}

The system prompt below is used in the probing setup of Chapter~\ref{chap:probing-results}, shared by all eight models. \texttt{\$\{collisions\}} is the same shared taxonomy listing from Section~\ref{sec:appendix-taxonomy}.

\begin{promptverbatim}
You are a coordinator that routes an incoming developer request to exactly
one specialist expert. Read the request, decide which single expert should
handle it, and route it by calling that expert's tool exactly once. The
experts are not topic areas: they look at the same system and optimise for
different constraints, so route on which constraint the developer is
actually held to, not on which one they talk about most.

Pairs that pull against each other:
${collisions}
\end{promptverbatim}

\section{Sub-Agent Tool Definitions}
\label{sec:appendix-subagent-tools}

Section~\ref{sec:ch7-probing-setup} states that each sub-agent is declared to the coordinator as a function-calling tool, following the schema of Section~\ref{sec:structured-output}. Below are the five tool definitions, one per sub-agent, exactly as each is rendered into the prompt alongside the system prompt of Section~\ref{sec:appendix-coordinator-prompt}. The \texttt{name} field is the sub-agent's identifier, and the \texttt{description} field states its role, what it handles, and an example request that routes to it. The \texttt{parameters} field defines an empty argument schema: the coordinator names a tool by calling it, and supplies no further argument.

\begin{promptverbatim}
{
  "type": "function",
  "function": {
    "name": "time_to_market",
    "description": "Time-to-Market expert. Handles: Absolute time-to-market. Accepts hardcoding, skipped abstractions and technical debt to get something demonstrable now.. Example request that routes here: \"we're demoing the ingest flow to the buyer on thursday and the deal is contingent on it running end to end. i'd rather wire the mapping properly but that can wait.\"",
    "parameters": {
      "type": "object",
      "properties": {}
    }
  }
}
\end{promptverbatim}

\begin{promptverbatim}
{
  "type": "function",
  "function": {
    "name": "scale_arch",
    "description": "Scalable Architecture expert. Handles: Future concurrency and horizontal growth. Decouples components, removes shared state, designs for many times the current load.. Example request that routes here: \"product has committed us to onboarding the three biggest tenants next quarter, which is roughly twenty times the events we take now. our worker pulls from one queue and holds a db connection open the whole time.\"",
    "parameters": {
      "type": "object",
      "properties": {}
    }
  }
}
\end{promptverbatim}

\begin{promptverbatim}
{
  "type": "function",
  "function": {
    "name": "resource_opt",
    "description": "Resource Optimization expert. Handles: Hard CPU, memory or cost ceilings. Strips abstractions and takes low-level control to fit inside a fixed budget.. Example request that routes here: \"finance capped our per-node spend and we can't move off the 2gb instances until next fiscal year. the parser copies every frame into a fresh buffer before decoding it.\"",
    "parameters": {
      "type": "object",
      "properties": {}
    }
  }
}
\end{promptverbatim}

\begin{promptverbatim}
{
  "type": "function",
  "function": {
    "name": "legacy_compat",
    "description": "Legacy Compatibility expert. Handles: Strict backward compatibility. Bound by published API contracts, existing clients and outdated dependencies that cannot move.. Example request that routes here: \"two external customers still call the v1 export endpoint and their contracts say we can't change its response shape for another year. i want to rework what happens behind it.\"",
    "parameters": {
      "type": "object",
      "properties": {}
    }
  }
}
\end{promptverbatim}

\begin{promptverbatim}
{
  "type": "function",
  "function": {
    "name": "strict_compliance",
    "description": "Strict Compliance expert. Handles: Security, privacy and regulatory standards. Zero-trust data handling, auditability, and rules that cannot be traded away.. Example request that routes here: \"our dpo signed off on a retention policy that says raw uploads never leave the audited bucket. right now the fine-tuning job reads straight from it by filename prefix.\"",
    "parameters": {
      "type": "object",
      "properties": {}
    }
  }
}
\end{promptverbatim}
